% ════════════
% mdi_prospectus.tex
% Material Dignity Infrastructure: A Financial Specification for Eliminating Chronic Urban Displacement
% Material Dignity Infrastructure Series
% ════════════

\documentclass[12pt, letterpaper]{article}

% ── PACKAGES ──
\usepackage[left=0.7in, right=0.7in, top=1in, bottom=1in, headheight=14pt]{geometry}
\usepackage[expansion=false]{microtype}
\usepackage{textcomp}
\usepackage{setspace}
\usepackage{booktabs}
\usepackage{tabularx}
\usepackage[titles]{tocloft}
\usepackage{tabularray}
\usepackage{longtable}
\usepackage{array}
\usepackage{multirow}
\usepackage{hyperref}
% natbib not used — manual thebibliography
\usepackage{amsmath}
\usepackage{amssymb}
\usepackage{parskip}
\usepackage{titlesec}
\usepackage{fancyhdr}
\usepackage[table]{xcolor}
\usepackage{caption}
\usepackage{footnote}
\usepackage{enumitem}
\usepackage{tikz}
\usepackage{needspace}
\usepackage{etoolbox}
\usepackage{float}

% ── COLOR SYSTEM ───
\definecolor{auditblue}{HTML}{003366}
\definecolor{rowgray}{HTML}{F5F5F5}

% ── GLOBAL SPACING ───
\setstretch{1.4}
\setlength{\parindent}{0pt}
\setlength{\parskip}{8pt}
\hyphenpenalty=1000

% ── SECTION FORMATTING. ZERO BOLDING IN TEXT. ───
\titleformat{\section}
  {\normalfont\large\color{auditblue}}{\thesection.}{0.5em}{}
\titleformat{\subsection}
  {\normalfont\normalsize\color{auditblue}}{\thesubsection.}{0.5em}{}
\titlespacing*{\section}{0pt}{18pt}{6pt}
\titlespacing*{\subsection}{0pt}{12pt}{4pt}

% ── HEADER / FOOTER ───
\pagestyle{fancy}
\fancyhf{}
\rhead{\small\textit{{Material Dignity Infrastructure Series}}}
\lhead{\small\textit{Working Paper, June 2026}}
\cfoot{\thepage}
\renewcommand{\headrulewidth}{0.4pt}

% ── TABLE SETTINGS ───
\renewcommand{\arraystretch}{1.3}

% ── ORPHAN / WIDOW CONTROL ───
\clubpenalty=10000
\widowpenalty=10000
\displaywidowpenalty=10000

% ── BIBLIOGRAPHY ORPHAN PREVENTION ───
\AtBeginEnvironment{thebibliography}{\interlinepenalty=10000}

% ── HYPERREF CONFIGURATION ───
\hypersetup{
  colorlinks=true,
  linkcolor=black,
  citecolor=black,
  urlcolor=black,
  pdftitle={{Material Dignity Infrastructure}: {A Financial Specification for Eliminating Chronic Urban Displacement}},
  pdfauthor={{Charles J. DiBella}},
  pdfsubject={Working Paper},
  pdfkeywords={{Blended Finance, CalAIM, Institutional Capital, Material Dignity Infrastructure, Urban Economics}}
}
\setcounter{tocdepth}{2}

\begin{document}
\captionsetup[table]{labelfont=normalfont, labelsep=period, skip=6pt}

% ── FRONT MATTER ───
\begin{titlepage}
\begin{center}
\setstretch{1.4}
{\LARGE \textbf{Material Dignity Infrastructure}}\\[6pt]
{\large A Financial Specification for Eliminating Chronic Urban Displacement}\\[0.5cm]
{\normalsize \textbf{Charles J. DiBella}}\\[0.01cm]
{\normalsize Systems Architect}\\[0.01cm]
{\normalsize Working Paper, June 2026}\\[0.3cm]
\textbf{ABSTRACT}
\vspace{0.2cm}
\end{center}
\begin{minipage}{0.95\textwidth}
\setstretch{1.15}
\noindent

This specification proposes a 220 million USD blended capital architecture (40\% catalytic equity, 60\% senior debt) to solve chronic urban displacement via the Material Dignity Infrastructure (MDI) framework. By acquiring distressed commercial real estate beneath replacement cost and deploying relational stabilization protocols, the model generates permanent operational revenue through California Advancing and Innovating Medi-Cal (CalAIM) billing. The Tripartite Coalition governance structure legally isolates physical assets within a Delaware Statutory Trust, transferring operational execution to a Public Benefit Corporation. The resulting prototype converts sunk municipal emergency costs into collateralized, cash-flowing asset equity, eliminating first-mover demonstration risk and unlocking autonomous institutional replication.

\end{minipage}
\end{titlepage}

\pagenumbering{arabic}

% ── DOCUMENT METADATA ────
\newpage
\section*{Document Metadata}

\noindent\textbf{Keywords}\\[1pt]
Blended Finance, CalAIM, Institutional Capital, Material Dignity Infrastructure, Urban Economics

\vspace{8pt}
\noindent\textbf{JEL Classification}
\begin{itemize}[nosep, before={\vspace{-8pt}}, leftmargin=2em]
    \item G11 Portfolio Choice • Investment Decisions
    \item G23 Non-bank Financial Institutions • Financial Instruments • Institutional Investors
    \item H75 State and Local Government: Health • Education • Welfare • Public Pensions
    \item I18 Health: Government Policy • Regulation • Public Health
    \item R31 Housing Markets • Housing Supply and Demand
\end{itemize}

\vspace{8pt}
\noindent\textbf{Target eJournals}
\begin{itemize}[nosep, before={\vspace{-8pt}}, leftmargin=2em]
    \item Urban Economics \& Regional Studies eJournal
    \item Public Finance \& Municipal Bonds eJournal
    \item Social Enterprise \& Impact Investing eJournal
\end{itemize}

\vspace{8pt}
\noindent\textbf{Author's Note}\\[1pt]
This paper is the fifth working paper in the Material Dignity Infrastructure series. This dossier provides the required financial specification to convert theoretical homelessness resolution frameworks into investable real estate assets, directly solving the first-mover capital risk gap.

\vspace{8pt}
\noindent\textbf{Suggested Citation}\\[1pt]
DiBella, C.J. (2026). Material Dignity Infrastructure: A Financial Specification for Eliminating Chronic Urban Displacement. Material Dignity Infrastructure Series, June 2026.

\newpage
% ── TABLE OF CONTENTS ────
\vspace{12pt}
\begingroup
\setstretch{1.35}
\setlength{\cftbeforesecskip}{2pt}
\setlength{\cftbeforesubsecskip}{-2pt}
\tableofcontents
\endgroup
% ── BODY ────
\setstretch{1.35}

\section{EXECUTIVE SUMMARY}
Chronic urban displacement represents a structural architectural failure requiring systemic capitalization, rather than a moral deficit requiring philanthropic mitigation. Los Angeles currently absorbs two billion dollars annually in sunk operational costs responding to emergency biological decay. The Material Dignity Infrastructure (MDI) framework converts these catastrophic municipal liabilities into collateralized, cash-flowing real estate equity. 

We propose acquiring One California Plaza at deeply distressed valuations (sub-replacement cost) to establish the first fully functional MDI prototype. The facility will stabilize one thousand metabolically compromised individuals, generating fifty thousand dollars per capita in verified avoided municipal costs. Autonomous replication across global capital markets initiates immediately following the Day 365 performance audit. 

\clearpage
\section{CAPITAL STACK AND YIELD MECHANICS}
Executing the One California Plaza prototype requires 220,000,000 USD in total capital. This figure encompasses acquisition costs, facility adaptation, Day One stabilization reserves, and thirty-six months of capitalized interest. The capitalization structure explicitly bridges the gap between unproven operational theory and institutionally deployable financial instruments.

Tranche A (Catalytic First-Loss Equity) totals 88,000,000 USD (40\% of Capital Stack), targeting mission-driven family offices utilizing Program-Related Investments (PRIs). It carries a deferred four percent target yield. This tranche absorbs initial operational volatility during the facility stabilization phase, de-risking the senior debt.

Tranche B (Senior Concessionary Debt) totals 132,000,000 USD (60\% of Capital Stack), targeting Community Development Financial Institutions (CDFIs) and mission-aligned commercial lenders. It carries an eight percent current target yield. Tranche B operates as traditional commercial real estate debt, strictly secured against the physical asset.

\section{LEGAL ARCHITECTURE AND GOVERNANCE}
The MDI framework explicitly rejects municipal procurement dependency. The Tripartite Coalition establishes autonomous operational governance utilizing two distinct corporate entities to isolate liability from physical assets.

Entity One is a Delaware Statutory Trust (DST) serving as absolute title holder for the physical real estate. It employs zero operational staff. This absolute separation guarantees that operational liabilities cannot attach to the physical real estate. 

Entity Two is a California Public Benefit Corporation (PBC) serving as exclusive operational manager and clinical service coordinator. The PBC leases the facility from the DST. The PBC absorbs all operational liability. In the event of catastrophic operational failure, the PBC can be dissolved without jeopardizing the physical real estate held by the DST.

\clearpage
\section{SOVEREIGN REVENUE MECHANICS}
Operational sustainability for the MDI prototype relies entirely upon sovereign state funding via Medicare and Medicaid (Medi-Cal). The California Public Benefit Corporation bills managed care plans directly for services rendered to the stabilized population under the CalAIM Community Supports framework.

Core HCPCS Billing Codes include H0043 U6 (Housing Transition Navigation Services), H0044 U2 (Housing Deposits), T2050 U6 (Housing Tenancy and Sustaining Services), and T2033 U6 (Recuperative Care). 

Los Angeles County currently expends approximately 50,000 USD annually per unsheltered individual. The MDI prototype captures this value through CalAIM billing. Sovereign billing initiates immediately upon resident intake. Full operational break-even is projected at Day 270, with stabilized positive cash flow generated exclusively through sovereign reimbursement channels by Day 365.

\section{STRUCTURAL RISK MITIGATION}
Institutional capital strictly requires comprehensive risk isolation. Tranche A explicitly absorbs first-mover demonstration risk. Once Day 365 performance metrics validate CalAIM revenue streams, the demonstration risk reaches zero, unlocking traditional institutional refinancing.

Operational liability contagion is mitigated through the Tripartite Coalition, which legally isolates the physical asset within the Delaware Statutory Trust. Political interference is mitigated because the MDI prototype operates entirely independently of municipal funding. Revenue volatility is mitigated by a fully funded 10.56 million USD capitalized interest reserve, guaranteeing thirty-six months of Tranche B debt service regardless of operational billing velocity.

\end{document}
